\section{LTI-Systeme in der realen Welt}
\label{sec:irl_lti}
Die Eigenschaften von LTI-Systemen sind in der digitalen Audiobearbeitung besonders dann interessant, wenn Realweltsysteme als LTI-System interpretiert
und durch ihre Impulsantwort approximiert werden können.
Ein bekannter Anwendungsfall ist der Convolution Reverb, der das Nachhallverhalten eines Raumes durch Convolution eines Signals mit einer im Raum
gemessenen Impulsantwort nachstellt. Ddazu wird Akustik eines Raumes durch einen Schallimpuls angeregt, mit einem Mikrofon aufgenommen, und digital gespeichert. 
Aus dem Messergebnis wird die Impulsantwort des Raums berechnet.
Durch eine Convolution der Impulsantwort mit einem beliebigen Eingangssignalwird wird das Ergebnis der Anwendung der Transferfunktion des Raums
auf das Eingangssignal approximiert, was zu einem Halleffekt führt, der dem im gemessenen Raum vorhandenen Nachhallverhalten entspricht.
Der Raum selbst weist lineare Eigenschaften auf, die eingesetzten Schallwandler und
Signalverstärker verursachen jedoch leichte nonlineare Verzerrungen. Die auf die Schallwandlung folgende Ausbreitung durch 
die Raumluft ist widerum nicht vollständig Zeit-Invariant \cite{farina2000,farina2007a,prawda2024}. 
Durch geeignete Messverfahren in Kombination mit entsprechenden Verfahren zur Berechnung der Impulsantwort (\textit{Deconvolution}) können auch Systeme mit leicht
nonlinearen und Zeit-Varianten Eigenschaften vollständig charakterisiert werden \cite{farina2000,farina2007a}. 
Neben dem gut nachvollziehbaren Anwendungsfall des Convolution Reverbs sind Verfahren, welche Impulsantworten und Convolution nutzen, in der
digitalen Signalverarbeitung allgegenwärtig. Sie kommen auch bei der qualitativen Messung von Audio-Equipment, etwa der Frequenzgangbestimmung
zum Einsatz. Ebenso werden Impulsantworten von Lautsprechersystemen in der Simulation von Gitarrenverstärkern verwendet, um den Einfluss
verschiedener an Verstärker angeschlossener Lautsprecherkabinette auf das Signal zu simulieren. 